\section{Verification of the bounded proof constructors}
\label{sec:annex-p05-pa-bounded-proof-builders}

\begin{annexitem}{annex:p5.bounded-logical-templates}{thm:p05-bounded-logical-templates}{The new logical rows and their exact open lists}
\begin{proof}
Work in \(\PA\), using the particular graph equations of
\cref{annex:p5.checker-description}. Every induction below is induction
on an arithmetic formula, retaining its other numerical parameters. A
function application denotes its unique specified graph output. The joining
invariant of \cref{thm:p05-builder-row-templates} gives successful old rows,
one parent for every copied old root, and zero parents for the new last row.
It remains, at each use, to verify the appended row and its actual open list.

The two projections have tag 2 or 3, one premise, no parameters, and root
equal to the selected child of the conjunctive premise. Their open list is
that premise's list. The two injections have tag 4 or 5, one premise and
no parameters; the premise is the corresponding child of their disjunction
root. They retain its open list. The case operation has tag 6, premises in
the order major, left minor, right minor, and parameters \([j,k]\).
If the major root is \(c\mathbin{|}d\), the two minor roots must be the
same output formula. It filters label \(j\), formula \(c\), in the
left minor, and label \(k\), formula \(d\), in the right minor. Each
filter must decrease length. Its result is the major list followed by the
filtered left and right lists. These are the actual tag-6 equations.

Bottom elimination is tag 11 with one bottom premise and no parameters;
it retains that list. Reductio is tag 12 with one bottom premise and
parameter \([k]\); at conclusion \(c\) it removes the leaves with
marker \(Sk\) and formula \(\mathsf{Neg}(c)\), nonvacuously.
All these rows have marker zero, as stipulated by \(\mathsf J\).
The common row tests follow from the input formula types, the stated
premise arities and the earlier-reference property of the join. In
\(\mathsf{Keep}(H,R)\), conjunction introduction first concatenates the
shifted open lists of \(H,R\); the subsequent right projection retains
that entire list, although its conclusion is only \(f(R)\).

For the excluded-middle expression put \(d=c\mathbin{|}\mathsf{Neg}(c)\).
In the first contradiction the marked \(c\) leaf gives \(d\) by left
injection. Its other premise is \(\mathsf{Neg}(d)\), marker 2, so
tag 10 gives bottom. Tag 9 with parameter \([0]\) removes the one
\(c\) leaf, marker 1, and concludes \(\mathsf{Neg}(c)\). Its open
list now consists of the one \(\mathsf{Neg}(d)\) leaf. Right injection
gives \(d\); contradiction with a new copy of
\(\mathsf L(\mathsf{Neg}(d),2)\) gives bottom with two open
\(\mathsf{Neg}(d)\) leaves. Tag 12, conclusion \(d\), parameter
\([1]\), removes both. Both discharges are nonvacuous; the last open
list is empty. The join invariant gives the required parent counts even
though the same formula and marker occur at two distinct assumption rows.

This checks all additional templates and \(\mathsf{EM}\) internally.
The empty open list makes \(\mathsf{EM}\) safe for its formula root.
The packed MP identity in the body follows by substituting the two raw
lists in the actual success expression of \(\mathsf{mp}_{19}\): its
first block is the antecedent list and its second the shifted implication
list. Thus the argument reversal in that identity is literal, not an
exchange rule at a row.
\end{proof}
\end{annexitem}

\begin{annexitem}{annex:p5.bounded-sum-exhaustion}{thm:p05-bounded-sum-exhaustion}{Two numeral-sum exhaustion recursions inside PA}
\begin{proof}
We first establish the syntax facts used by the recurrence. Induction on
\(n\), using the syntax-table equations and the numeral invariant, gives
\(\mathsf{Fm}(\mathcal D(t,n))=1\) for a term \(t\). Its free and
all-variable indices are contained in those of \(t\). At zero its code
is bottom. At a successor its two children are the old disjunction and
the equality between \(t\) and a variable-free numeral. If \(t=v_x\),
its substitutions by variable-free numerals, and by a variable at another
index, are obtained by the same recursion on these two children; there
are no binders in \(\mathcal D\). The no-conflict theorem identifies
these computations with the selected substitution machine. In particular
the code with counter \(Sn\) ends in exactly the equality to \(\nu(n)\).

Retain pairwise distinct \(x,z,w\geq3\). The already verified axiom
instance chains give, for \(t=v_x,u=v_w\),
\[
\begin{array}{c|c|c}
 \epsilon&f(Z_\epsilon(t))&f(S_\epsilon(t,u))\\\hline
 0&\mathsf{Eq}(a(t,o),t)&\mathsf{Eq}(a(t,s(u)),s(a(t,u)))\\
 1&\mathsf{Eq}(a(o,t),t)&\mathsf{Eq}(a(s(u),t),s(a(u,t))).
\end{array}
\]
For the second orientation these are obtained by copying \(D_Z,D_L\)
and performing their specified universal eliminations. The replacements
at the outer index 0 contain no variable 1; their remaining matrices
after eliminating index 1 are quantifier-free. Thus the actual
substitution-machine roots are precisely the displayed codes. All their
open leaves are unmarked members of \(\mathcal A_*\). The calls to
\(\mathsf A_1,\mathsf A_2,\mathsf A_3\) below likewise use terms
with indices only among \(x,z,w\), hence none with index 1. In
\(\mathsf A_2\), the second replacement is a numeral; in
\(\mathsf A_3(v_z)\) the existential index 1 is retained. This verifies
every axiom-instance side condition used by the construction.

We prove the claimed safe root by induction on \(n\). For the base,
the stage's recursive argument is zero and only the \(n=0\) branch
is used. For the step at \(Sn\), its recursive argument is the
previous safe tree, with the exact universally quantified root in the
statement. The following calculation treats the common portions of the
base and step at a stage counter \(r\).

The outer marked assumption \(H\) has formula
\(g_\epsilon(v_x,v_z)=\nu(r)\), marker 1. The zero minor's
assumption \(Z\) has formula \(v_z=o\), marker 2. The zero minor
chains exactly
\[
 v_x=g_\epsilon(v_x,o),\qquad
 g_\epsilon(v_x,o)=g_\epsilon(v_x,v_z),\qquad
 g_\epsilon(v_x,v_z)=\nu(r).
\tag{1}
\]
The first equality is the symmetry of \(Z_\epsilon(v_x)\); the
second is the symmetry of congruence applied to \(Z\); the third is
\(H\). Both transitivity joins therefore have identical middle terms.
Right disjunction introduction concludes \(\mathcal D(v_x,Sr)\).
This branch retains exactly one copy of \(H\) and one copy of \(Z\),
besides unmarked allowed axiom leaves.

In the other minor, the predecessor assumption \(K\) is
\(v_z=s(v_w)\), marker 4. The definition of \(W\) chains
\[
 s(g_\epsilon(v_x,v_w))=g_\epsilon(v_x,s(v_w)),\qquad
 g_\epsilon(v_x,s(v_w))=g_\epsilon(v_x,v_z),\qquad
 g_\epsilon(v_x,v_z)=\nu(r).
\tag{2}
\]
The first equality is the symmetry of the declared successor instance,
and the second is the symmetry of congruence applied to \(K\).
Consequently its root is exactly
\(s(g_\epsilon(v_x,v_w))=\nu(r)\). Its marked leaves are one
\(H\) and one \(K\).

At \(r=0\), the instance of \(q_1\) at
\(g_\epsilon(v_x,v_w)\) negates precisely that equality. The
contradiction row gives bottom, and the declared bottom-elimination row
gives \(\mathcal D(v_x,1)\). Thus the base does not need a preceding
exhaustion proof or a predecessor of zero.

At \(r=Sn\), its numeral is literally \(s(\nu(n))\). The
\(q_2\) instance in the stage has antecedent equal to the root in
(2) and consequent
\(g_\epsilon(v_x,v_w)=\nu(n)\). Raw implication elimination takes
(2) first and this implication second. The induction hypothesis's proof
has root
\[
 \forall x_x\forall x_z\,
       (g_\epsilon(v_x,v_z)=\nu(n)\to\mathcal D(v_x,Sn)).
\]
Its first elimination uses \(v_x\) itself, which differs from the
remaining binder \(z\); its second uses \(v_w\) in a
quantifier-free matrix. Therefore the two selected substitutions give
literally
\(g_\epsilon(v_x,v_w)=\nu(n)\to\mathcal D(v_x,Sn)\).
The next raw MP gives \(\mathcal D(v_x,Sn)\). Left disjunction
introduction appends equality to \(\nu(Sn)\), yielding exactly
\(\mathcal D(v_x,S(Sn))\), the required stage root. All copied
induction-hypothesis leaves are unmarked allowed axioms; \(H,K\)
still survive once in this minor.

To discharge the predecessor minor, the major is the actual \(q_3\)
instance applied to \(\neg(v_z=o)\), marker 3. Its root is
\(\mathsf{Ex}(1,\mathsf{Eq}(v_z,s(v_1)))\). The tag-16
parameters are exactly \([1,\mathsf{Eq}(v_z,s(v_1)),w,3]\).
The selected matrix substitution gives the formula of \(K\), and
the label-3 filter removes its one marker-4 occurrence. Nonvacuity holds.
The conclusion's all-variable set is contained in \(\{x\}\),
the major formula uses only \(z,1\), and the remaining temporary
assumptions use only \(x,z\). Every remaining axiom leaf uses only
indices 0 and 1. As \(w\geq3\) differs from \(x,z\), each actual
\(\mathsf{AV}\) factor is zero at \(w\), so the full freshness
product is one. The open-list order is the minor remainder, then the
major list. It retains one \(H\) and the marker-3 negative-zero
assumption.

The final tag-6 row uses the empty-context tree \(\mathsf{EM}(e)\)
as major. It removes marker 2, formula \(e\), in the zero minor,
and marker 3, formula \(\mathsf{Neg}(e)\), in the other minor.
Both occur once. Its remaining list includes two copies of \(H\),
one from each minor, together with unmarked allowed axiom leaves. Tag 7,
parameter \([0]\), removes both copies of the exact formula \(a_r\).
This is nonvacuous. Both universal introductions now have only closed
axiom leaves as assumptions, so their free-variable products are one.
Their order, first \(z\) and then \(x\), gives the exact displayed
root. No assumption with marker 1, 2, 3 or 4 survives.

Each calculation above is an identity and local-rule check for the actual
row expressions, so the accepted join theorem establishes all successful
row bits and all parent counts. Its shifted-leaf equations prove the leaf
assertions with their multiplicities. This is a base and successor proof
of one arithmetic safe-tree invariant, universally in the retained index
parameters. Applying \cref{cor:p05-pa-recursion-invariant} proves it
for all numerical \(n\). The guard excludes invalid triples throughout
and has the stated zero default. The final formula has no free indices,
by the syntax-table equations, so packing gives the selected sentence
proof. This argument verifies both explicitly listed orientations.
\end{proof}
\end{annexitem}

\begin{annexitem}{annex:p5.bounded-range-exhaustion}{thm:p05-bounded-range-exhaustion}{Eliminating the comparison witness and retaining its binder}
\begin{proof}
First work on the triple domain of \(\mathsf{CS}\). Its marked
assumption is \(v_x+s(v_z)=\nu(n)\). Symmetry of the actual
\(q_5\) instance, followed by transitivity with that assumption,
gives exactly
\[
                       s(a(v_x,v_z))=\nu(n).
\tag{3}
\]
At zero, \(q_1\) at \(a(v_x,v_z)\) negates (3), so the declared
contradiction has root bottom, which is \(\mathcal D(v_x,0)\).
At \(n=Sr\), the \(q_2\) instance peels (3) to
\(a(v_x,v_z)=\nu(r)\). The two universal eliminations from
\(\mathsf C_0(x,z,w,r)\), first at \(v_x\), then at \(v_z\),
have exact identity substitutions and root
\(a(v_x,v_z)=\nu(r)\to\mathcal D(v_x,Sr)\).
The specified raw MP gives the required disjunction. In both cases the
one original marked assumption survives; tag 7 removes it, and the two
universal introductions close \(z\) then \(x\). Only unmarked
\(\mathcal A_*\) leaves remain. The axiom-instance substitutions
have indices at least 3 and satisfy the same no-conflict checks as in
the preceding proof. Hence every row test, parent count and safe-root
assertion for \(\mathsf{CS}\) follows internally.

Consider now \(\mathsf E_\diamond(i,h,n)\), with \(i\ne h\).
Its fresh choices satisfy \(y,x,z,w>\max(2,i,h)\) and are
pairwise distinct. In the weak case the bound assumption and its minor
have exact formulas
\[
 \exists x_h(x_h+x_i=\bar n),\qquad x_y+x_i=\bar n.
\tag{4}
\]
In the strict case they have exact formulas
\[
 \exists x_h(x_i+Sx_h=\bar n),\qquad x_i+Sx_y=\bar n.
\tag{5}
\]
The comparison matrix has no binder, so substitution at \(h\) gives
the displayed minor without any renaming. The weak case's \(Q\)
is the second sum orientation, which enumerates the right addend;
its two eliminations at \(v_i,v_y\) have root
\(x_y+x_i=\bar n\to\mathcal D(v_i,Sn)\). The strict case's
\(Q\) is \(\mathsf{CS}\), giving
\(x_i+Sx_y=\bar n\to\mathcal D(v_i,n)\).
At the first substitution the remaining binder \(z\) is fresh from
\(v_i\); at the second the matrix is quantifier-free. Thus these
are literal selected-machine identities even when \(i=0\) or
\(i=1\). No arithmetic axiom with an unsafe outer replacement by
\(v_1\) is being instantiated at these calls: the already closed
exhaustion theorem has the fresh outer binders just specified.

The MP tree uses the marked minor first and the instantiated implication
second. Its sole marked open leaf is the minor, marker 2. The tag-16
parameters \([h,c,y,1]\) remove that leaf nonvacuously. The conclusion
has only the possible index \(i\), the major has only \(i,h\), and
the remaining open leaves are the original bound assumption and unmarked
\(\mathcal A_*\) leaves. Hence \(y>\max(2,i,h)\) makes all three
freshness factors and the open-leaf product one. The result retains the
original bound assumption, marker 1, once. Implication introduction with
label 0 removes it and gives exactly the announced exhaustion implication.
The comparison binder \(h\) occurs in that antecedent exactly as in
(4) or (5); no alpha conversion was used.

At the strict zero bound, (3) supplies the contradiction and the last
exhaustion implication concludes bottom. At the weak zero bound,
\(\mathcal D(v_i,1)=\mathsf{bc}\mathbin{|}\mathsf{Eq}(v_i,o)\);
the one possible value remains present. The earlier join theorem gives
every parent count and the full selected row invariant in both cases.
\end{proof}
\end{annexitem}

\begin{annexitem}{annex:p5.bounded-case-elimination}{thm:p05-bounded-case-elimination}{Equality transport, finite cases and universal introduction}
\begin{proof}
We need an exact substitution identity for arbitrary formula-code input.
For \(\mathsf{Fm}(c)=1\), the no-conflict table gives
\[
                    \mathsf{sub}(c,i,v_i)=c.
\tag{6}
\]
Prove it by the table's strong induction on expression code. At a variable,
replacing \(v_i\) by itself changes no code. Zero and bottom are
unchanged. Each nonbinding constructor commutes with the child identities.
At a binder for \(i\), substitution stops and returns the entire node.
At another binder \(j\), the only variable of the replacement is
\(i\ne j\), so there is no conflict and the matrix induction instance
returns the original matrix. These cases prove success of the table and
its unchanged output; \cref{thm:p05-builder-nocapture} identifies it
with the selected machine. This includes shadowing of index \(i\)
and occurrences that are only bound. The replacement is a valid term,
as required by that theorem.

Fix \(c,i,Q\), with \(c\) a formula. Induct on \(n\) under
the premise that each earlier entry \(k<n\) is safe with root
\(\mathsf{sub}(c,i,\nu(k))\). The invariant for
\(\mathsf{Cases}(c,i,Q,n)\) is: successful tree, root \(c\),
exactly one marked open leaf with formula \(\mathcal D(v_i,n)\)
and marker 1, and all other open leaves unmarked members of
\(\mathcal A_*\). The statement about one marked leaf means one
entry in the actual open-index list; its formula and marker are row
projections. It is consequently an arithmetic formula for induction.

At zero, bottom elimination from the one marked bottom leaf has root
\(c\). Since \(\mathcal D(v_i,0)=\mathsf{bc}\), the invariant
holds directly, with no unmarked leaves. In the step the new right minor
starts with \(v_i=\nu(n)\), marker 1. The exact symmetry template
gives \(\nu(n)=v_i\). Its second equality-elimination premise is
the supplied safe tree \(e(Q,n)\), with root
\(\mathsf{sub}(c,i,\nu(n))\). Tag 19 has precisely parameters
\([\nu(n),v_i,c,i]\); by (6) its conclusion is \(c\). Its only
marked open leaf is the new equality assumption, and its remaining leaves
are those of the supplied case proof.

The tag-6 major is the new marked
\(\mathcal D(v_i,Sn)=\mathcal D(v_i,n)\mathbin{|}
\mathsf{Eq}(v_i,\nu(n))\) leaf. The left minor is the preceding
state and the right minor is the just-described transport. Their roots
are both \(c\). Parameters \([0,0]\) remove the left minor's
unique \(\mathcal D(v_i,n)\) leaf and the right minor's unique
equality leaf. Both filters decrease length; each sees only its own
premise branch. It therefore causes no extra removal that the label
also equals the label of the major. The major remains once in the new
open list, exactly as required. The branch shift and parent-count
equations of the join theorem prove every old and new row assertion.
This completes the internal induction.

The tag-7 row defining \(\mathsf{CI}\) now removes the one marked
\(\mathcal D(v_i,n)\) leaf and has exact root
\(\mathcal D(v_i,n)\Rightarrow c\). It retains only unmarked
allowed axiom leaves, proving the first assertion. No requirement that
\(x_i\) occur freely in \(c\) was needed: equality elimination
allows an unused matrix variable, and the new equality assumption is
still used as its first premise. This is why the discharge remains
nonvacuous even for a constant matrix.

For the quantified constructor retain its additional premises. The hole
\(j=1+\max(i,h)\) differs from both \(i,h\). The term \(u\)
and the numeral \(\nu(b)\) have no variables. Applying the selected
no-conflict equations to the guard context gives exactly
\[
\begin{aligned}
 \mathsf{sub}(G_\diamond(i,h,v_j),j,u)&=G_\diamond(i,h,u),\\
 \mathsf{sub}(G_\diamond(i,h,v_j),j,\nu(b))
                                      &=G_\diamond(i,h,\nu(b)).
\end{aligned}
\tag{7}
\]
The existential comparison binder remains \(h\). Thus the first
tag-19 row uses the equality proof \(R_u\) first and the original
marked guard second; its root is the normalized guard in the second
line of (7). Its marked open list still contains the original guard
once. Apply the exhaustion implication by the first raw MP, and the
safe case implication by the second. The final root is \(c\), and
the only marked open leaf is still the original \(G\), marker 1.
Tag 7 with parameter \([0]\) removes it nonvacuously. This gives
\(G_\diamond(i,h,u)\Rightarrow c\), with the original bound term
\(u\) in the antecedent. Its other open leaves are unmarked members
of \(\mathcal A_*\); all are sentences. Tag 15 with parameter
\([i]\) therefore passes its free-variable product and gives exactly
the original outside binder \(i\).

Since \(u\) has no variables, the guard's only free index is \(i\).
The premise concerning \(c\) gives the same containment for the
matrix. The outer quantifier removes that index. The syntax-table equations
therefore prove its sentence test, and the safe-tree assertion gives every
factor of the selected packed proof test. At \(N=0\), the case
recursion consists of its bottom-elimination base; the exhaustion proof
actually derives bottom from the guard. The same implication and universal
rows close it. Thus an empty list of instances has not been generalized
without a derivation.
\end{proof}
\end{annexitem}

\begin{annexitem}{annex:p5.fixed-bounded-builders}{thm:p05-fixed-bounded-builders}{Fixed formation, both signs and all bounded quantifier cases}
\begin{proof}
We establish finite describability, literal substitution, the Boolean
equivalence and the two raw checker invariants for the formula fixed in
the body. This is a metainduction on that one finite marked formation.
Every induction assertion delivered by a formation step is a theorem of
\(\PA\) with its numerical tuple universally quantified.

First verify the needed finite descriptions. The logical templates and
each stage of the exhaustion constructions are compositions of earlier
descriptions and finite conditionals. Their recursions use the last
counter and the supplied preceding state only. The two exhaustion
orientations have separate finite formations. \(\mathsf{CS}\)
calls only the already defined \(\mathsf C_0\). Comparison
exhaustion calls those fixed functions. The case recursion calls the
earlier list entry and equality templates. In the formula construction,
the Boolean and proof descriptions at a child are already fixed before
the parent is formed. A bounded sum, product, search or list fold calls
that child's description with specified tuple projections. Its bound is
the parent's already evaluated arithmetic term. All failure outputs and
branch priorities are specified. The graph theorem therefore supplies
total unique values for each of these descriptions. No stage evaluates
an arbitrary input description or searches for an arbitrary proof.

Here are the required substitution identities. Numerals have no variables,
and their substitution invariance has the replacement-term hypothesis
provided by \cref{thm:p05-builder-nocapture}. For distinct indices
\(i,j\), substitution by two numerals commutes literally. Strong
induction in the no-conflict table proves this: variables give the two
specified replacements or remain unchanged; nonbinding constructors use
their child instances. At a binder for \(i\), the \(i\)-operation
stops, while the \(j\)-operation commutes with the binder and changes
only the matrix's free \(j\)-occurrences; applying \(i\) after it
still stops at that same binder. The case of a binder for \(j\) is
the reverse. At every other binder both replacements are variable-free,
so neither operation renames it and the matrix induction instance
applies. The same argument covers terms.

Consequently any prescribed finite order of the free-index substitutions
has its literal simultaneous-numeral value. At a quantifier with current
binder \(i\), perform the substitutions for its other free indices
and call the resulting matrix \(c\). For every numerical \(k\),
the child description's numeral code is exactly
\[
                       \psi^\nu(\bar m,k)=\mathsf{sub}(c,i,\nu(k)).
\tag{8}
\]
Only the free occurrences of the current \(i\) are replaced; an
inner binder for \(i\) stops this operation. At a connective, the
corresponding child-code constructor equation follows from the same
table. At a marked quantifier, the full code is exactly its retained
binder around \(G\mathbin{\&}c\) or \(G\Rightarrow c\).
Its bound has no occurrence of the current bound variable; after the
other parameters are replaced it is the variable-free term \(u\).
The guard's comparison binder was fixed before these substitutions,
and none of these closed replacements renames it or chooses a new one.
This proves the necessary literal root equations, including parameters
that share a name with an inner shadowing binder.

The elementary numerical tests are zero-one and agree in \(\PA\)
with the corresponding comparisons of the fixed term values, by the
accepted fixed-term and literal theorem. Thus the atomic Boolean
equivalence holds internally. For later folds, induction on a counter
\(r\) gives the following internal facts for any separately fixed
zero-one expression \(b(k)\):
\[
\begin{aligned}
 \prod_{k<r}b(k)=1&\ \leftrightarrow\ \forall k<r\,b(k)=1,\\
 \mathsf{sg}\left(\sum_{k<r}b(k)\right)=1
                      &\ \leftrightarrow\ \exists k<r\,b(k)=1.
\end{aligned}
\tag{9}
\]
At zero the product is one and the sum zero, giving the empty universal
and existential cases. At a successor the product is one exactly when
both the preceding product and the next bit are one. The sum is positive
exactly when its preceding sum is positive or its next bit is one;
its terms are nonnegative. These arithmetic facts follow from addition
and multiplication's defining equations and order laws in
\cref{annex:p5.internal-pr-graphs}. They prove (9) by \(\PA\)
induction. The product is itself zero-one; when it is zero there is a
zero factor, by the first equivalence and classical quantifier negation.
When the sign of the sum is zero, every factor is zero, by the second.

The stated least-search fold selects an actual witness whenever the
corresponding existential in (9) holds. To verify this in \(\PA\),
induct on its counter: before the first success its flag is zero and
no earlier index succeeds; after a success its flag is one, its saved
index is below the counter, its test succeeds, and every smaller index
fails. The next update either records the current first success or
retains the saved data. At the supplied final bound this proves both
the witness property and leastness. Failure returns the declared zero
without any witness assertion. The list-fold induction similarly gives
length \(r\) and entry \(R_\psi^\pm(\bar m,k)\) at every \(k<r\),
since appending preserves the prefix and adds exactly that next entry.
These are arithmetic inductions for the specified folds.

At bottom the Boolean bit is zero. The positive implication is obtained
from its contradictory arithmetic antecedent; no successful zero raw
record is asserted. The negative raw record is the accepted explicit
two-row proof of \(\neg\bottom\). At equality, the tests select
the accepted equality or disequality proof and the two term transports.
At a weak comparison the success tests are \(u^\#\leq v^\#\)
and \(u^\#>v^\#\); at a strict comparison they are
\(u^\#<v^\#\) and \(u^\#\geq v^\#\). The six corresponding
literal root identities, the seven-axiom convention and the stronger
\(\mathcal A_*\) leaf assertion are exactly those proved in
\cref{thm:p05-numeral-literal-builders,thm:p05-fixed-term-literal-builders}.
Under a success premise the packed literal output has its literal
outer tags, so the declared decoder returns the very raw tree already
verified there. Comparison binders are passed through unchanged.

For negation, the positive tree is the child's negative tree, whose
root is already \(\mathsf{Neg}(c)\). In the negative case the child
positive tree has root \(c\). The new assumption
\(\mathsf{Neg}(c)\), marker 1, is the second premise of a
contradiction. Tag 9 removes it and concludes
\(\mathsf{Neg}(\mathsf{Neg}(c))\). It does not remove a negation
from the target. The arithmetic truth clauses follow from the child's
zero-one bit and classical negation.

For conjunction, bit one supplies both positive child invariants; tag 1
gives exactly \(c\mathbin{\&}d\) and concatenates their axiom
leaves. At bit zero, test the left bit first. If it is zero, project
the left conjunct from the marked conjunction assumption and contradict
the left negative proof. If it is one, the right bit must be zero;
project the right conjunct and contradict the right negative proof.
In each case the only marked leaf is the original conjunction assumption,
which survives once and is removed by tag 9. Its conclusion is exactly
the negation of the conjunction. The zero-one product equation gives
the internal Boolean equivalence in both directions.

For disjunction, bit one selects the left positive proof when its bit
is one; otherwise the right bit is one and the right positive proof
is selected. Tags 4 and 5 have exactly the required original disjunction
root. At bit zero both child bits are zero. Each minor contradiction
uses its marked child formula, marker 1, followed by the corresponding
negative child proof. Their roots are both bottom. The tag-6 major is
the marked original disjunction, also marker 1; its two branch-specific
filters, with parameters \([0,0]\), remove the marked child in
each minor and retain the major. This remains valid when the two child
formulas coincide, since the filters still consult separate branches.
The final tag-9 row removes the one major assumption and returns the
literal negation of the disjunction. Sign of the sum gives the
arithmetic equivalence, using nonnegativity and the two possible bits.

For implication, bit one has two cases. When the antecedent bit is
zero, the marked antecedent and its negative proof give bottom;
bottom elimination gives the consequent, retaining that antecedent
once. When its bit is one, the consequent bit is one. The operation
\(\mathsf{Keep}(\mathsf L(c,1),R_\chi^+)\) gives the consequent
and retains the marked antecedent once. Tag 7 removes it in either
case, nonvacuously, and returns the original implication. At bit zero
the antecedent bit is one and the consequent bit zero. Raw MP consumes
the positive antecedent proof before the marked implication, giving
the consequent; the negative consequent proof gives bottom. Tag 9
removes the implication assumption, which survives once. The stated
conditional's two numerical branches prove its Boolean equivalence
to the original implication. These arguments cover all eleven connective
branches and every success condition; the failure
branches are total but make no proof claim.

For a marked quantifier, the fixed-term theorem identifies its numerical
bound value \(b\) internally and supplies the safe equality tree
\(R_t\) with root \(u=\nu(b)\). The arithmetic comparison laws
identify its range with \(k<N_\diamond(b)\), retaining the
original \(<\) or \(\leq\) formula as the guard. Using (8), the
child Boolean equivalence and (9) proves the parent Boolean equivalence
for the marked existential and universal. This is a proof of the
fixed original arithmetic formulas from their graph equations; it
does not appeal to their truth in the standard structure.

In the positive existential case the sum-sign bit supplies the least
successful \(k<N\) through the verified search. Hence the actual
bound comparison holds and the child positive invariant holds at that
same numerical \(k\). The fixed literal constructor supplies
\(L_k\) with root \(\mathsf{sub}(G,i,\nu(k))\). The conjunction
tree has root
\[
 \mathsf{sub}(G,i,\nu(k))\mathbin{\&}\mathsf{sub}(c,i,\nu(k))
               =\mathsf{sub}(G\mathbin{\&}c,i,\nu(k)).
\tag{10}
\]
The no-conflict table proves this exact constructor equation. The
existential-introduction row has parameters \([i,C,\nu(k)]\) and
therefore concludes the original \(\mathsf{Ex}(i,C)\). It retains
only the two proofs' allowed unmarked axiom leaves. The witness is
the selected numerical value's numeral, never the numerical value as
a term code.

In the negative universal case the product bit zero supplies a least
\(k<N\) at which the child bit is zero. The bound proof \(L_k\)
again has its exact guard root. Universal elimination from the marked
original \(A\), with parameters \([i,G\Rightarrow c,\nu(k)]\),
gives the implication from that guard to the exact child instance.
Raw MP gives the instance; the negative child proof contradicts it.
The one marked \(A\) leaf survives once, and tag 9 removes it to
give exactly \(\mathsf{Neg}(A)\). These two witness constructions
require no eigenvariable and introduce no additional axiom leaves.

For the positive universal, the product bit one gives every child
positive invariant for \(k<N\). The verified list fold therefore
supplies all of the case premises for
\cref{thm:p05-bounded-case-elimination}. Equation (8) identifies their
roots with its exact required substitutions. The matrix \(c\) has
only the possible free index \(i\), since all its other free
parameters were replaced by numerals. That theorem gives a safe tree
with exact root \(\mathsf{All}(i,G\Rightarrow c)\). It uses
the disjunction-exhaustion implication, the case-elimination tree,
implication introduction and universal introduction specified in
the body. Thus the list fold is only the supply of minor-premise
proofs to a separately constructed quantified derivation.

For the negative existential, the sum-sign bit zero gives every child
negative invariant. The negative list is consequently safe with roots
\(\mathsf{Neg}(\mathsf{sub}(c,i,\nu(k)))\). This equals
\(\mathsf{sub}(\mathsf{Neg}(c),i,\nu(k))\) by the nonbinding
substitution equation. The same universal constructor gives \(U\)
with exact root \(\mathsf{All}(i,G\Rightarrow\mathsf{Neg}(c))\).
Its only open leaves are unmarked members of \(\mathcal A_*\).

The new eigenvariable \(y=1+\max(2,m(E))\) is absent from every
variable occurrence of the major \(E=\mathsf{Ex}(i,G\mathbin{\&}c)\),
including every original binder, and is at least 3. Its replacement
therefore passes every no-conflict test; it also differs from \(i\).
Writing \(G_y=\mathsf{sub}(G,i,v_y)\) and
\(c_y=\mathsf{sub}(c,i,v_y)\), the minor leaf has exact formula
\(G_y\mathbin{\&}c_y\). The two separate copies of this leaf
give \(G_y\) and \(c_y\) by the two projections. The elimination
from \(U\) has literal root \(G_y\Rightarrow\mathsf{Neg}(c_y)\);
no binder of \(G\) or \(c\) equals the fresh \(y\), so the
selected machine commutes without renaming. Raw MP gives
\(\mathsf{Neg}(c_y)\), and contradiction with \(c_y\) gives
bottom. Exactly two marker-2 copies of the minor survive. The tag-16
parameters are \([i,C,y,1]\), so the filter removes both copies of
the exact selected minor formula and is nonvacuous.

The conclusion bottom has no variables; the major has none at \(y\)
by its choice. After the minor discharge, every remaining leaf is an
unmarked allowed axiom or the one major assumption \(E\), marker 1.
The allowed axiom leaves use only indices 0 and 1. Hence every factor
of the tag-16 all-variable freshness test is one, including its test
on the remaining major assumption. Its open list is the filtered minor
list followed by the major list, as specified. The final tag-9 row
removes that sole marked major and concludes precisely
\(\mathsf{Neg}(E)\). This verifies the fourth bounded-quantifier
case in both comparison orientations.

At a strict bound of zero, \(N=0\). The existential bit is zero
and the universal bit one. The positive existential and negative
universal success guards fail; no least-search witness is asserted
there. The positive universal uses the case recursion's bottom base
and the strict-zero exhaustion proof. The negative existential uses
the same valid universal constructor for \(\neg c\) and the
two-copy minor just checked. A weak bound of zero has \(N=1\),
so its zero witness and zero counter case are included instead of
being discarded. These facts also cover an unused quantified variable:
every guard is consumed, and the case tree uses its equality assumption
even if the target matrix is invariant under that substitution.

For dependent bounds, all preceding arguments keep the outer numerical
tuple as parameters. At an inner quantifier the currently supplied
witnesses are already entries of that tuple; its own bound is evaluated
at those entries before its fold. Its binding variable is absent from
the bounding term, so substituting that value for the range parameter
does not change or capture a bound variable. Equation (8) continues
to identify each child proof's numeral tuple with the literal matrix
instance. For example, in a marked expression
\(\forall x_i<t(\bar z)\,\exists x_j<s(\bar z,x_i)\,\theta\),
the outer case at \(k\) calls the inner description with that
very \(k\); the inner search bound is
\(N_\diamond(s^\#(\bar m,k))\), and its returned witness \(l\)
is used in the term \(\nu(l)\). The child instance is the result
of the two specified substitutions by \(\nu(k),\nu(l)\), which
commute at distinct indices by the proved table identity. If the
inner binder reuses the outer index, substitution for the outer
occurrence stops there and the inner tuple projection overrides it
only in that inner scope. This is exactly the shadowing case proved
above, not a flattening of witnesses to one common bound.

We have proved both raw safe-tree assertions at every formation step.
Every output root after all free parameters receive numerals is a
sentence by the same syntax-table induction: the only newly free
index at a marked matrix is removed by its original outer binder.
Negation preserves its zero free-variable bits. Thus the safe-tree
assertions supply every actual factor of \(\chi_{19}\) after
packing: nonempty successful tree, required parent counts, sentence
root, literal root equality, and the axiom product over its exact
open list. Every member of \(\mathcal A_*\) was already verified
to pass \(\mathsf{Ax}_{19}\). The final success guards preserve
these packed values; their failure values remain zero.

Finally expand each fixed graph application through its unique output
variable, existentially introduce the constructed \(p\), and use
the Boolean equivalence to replace its sign condition by
\(\varphi(\bar m)\) or \(\neg\varphi(\bar m)\). Implication
introduction and universal closure give exactly the two formulas in
the body. Totality, defining equations, uniqueness and the explicit
base/step proofs supply these arithmetic derivations; correctness
of a standard computation alone was never used as a universal premise.
With no free parameters the tuple projections and prefix are empty,
and the same finite construction is nullary. This concludes the
fixed bounded-formula theorem at its stated scope.
\end{proof}
\end{annexitem}
